% SPDX-License-Identifier: CC-BY-SA-4.0
% Author: Matthieu Perrin
% Part: <Nom de la partie>
% Section: <Nom de la section>
% Sub-section: <Nom de la sous-section>  % (facultatif, laisser vide si non utilisé)
% Frame: <Titre de la slide>

\begingroup

\begin{frame}{Généralisation : le théorème de Rice}

  Soit $\langle \mathcal{L}, \llbracket \cdot \rrbracket \rangle$ un encodage effectif. 
  \begin{block}{Définition -- Problème sémantique non-trivial}
    \begin{itemize}
    \item Un problème de décision est dit \structure{sémantique}
      s'il est de la forme
      \Probleme{Sem$_S$}{
        Un programme \structure{$P \in \mathcal{L}$}
      }{
        Est-ce que \structure{$\llbracket P \rrbracket \in S$} ?
      }
      où $S \subseteq \textsc{re}$ est la \structure{propriété sémantique} associée
    \item On dit que $\textsc{Sem}_S$ est \structure{non-trivial} si \alert{$S \neq \emptyset$} et \alert{$S \neq \textsc{re}$}
    \end{itemize}

  \end{block}
  
  \begin{block}{Théorème -- Rice (1953)}
    Tout problème sémantique non-trivial est indécidable.
  \end{block}  

\end{frame}

\endgroup
